% Banco de ejercicios — capítulo 4
% El capítulo viene determinado por el nombre de este archivo.

\begin{exo}[Integrar una forma diferencial a lo largo de un camino]{integrale-forme-differentielle}
\keywords{forma diferencial, diferencial exacta, criterio de Schwarz, integral de línea}

Se consideran las dos formas diferenciales
\[\omega_1=y\,\mathrm{d}x+x\,\mathrm{d}y,\qquad \omega_2=y\,\mathrm{d}x,\]
y tres caminos que unen el origen $O(0,0)$ con el punto $M(1,1)$: $\gamma_1$,
formado por el segmento horizontal $O\to(1,0)$ seguido del vertical
$(1,0)\to M$; $\gamma_2$, formado por el segmento vertical $O\to(0,1)$
seguido del horizontal $(0,1)\to M$; y $\gamma_3$, el segmento rectilíneo que
une directamente $O$ con $M$.

\begin{enumerate}
\item Demostrar que $\omega_1$ es exacta hallando una función $f$ tal que $\omega_1=\mathrm{d}f$.
\item Calcular $\int_{\gamma_1}\omega_1$ y $\int_{\gamma_2}\omega_1$ parametrizando cada segmento. Obtener de nuevo el resultado sin calcular.
\item Aplicar el criterio de Schwarz a $\omega_2$. ¿Qué se puede concluir?
\item Calcular $\int_{\gamma_1}\omega_2$, $\int_{\gamma_2}\omega_2$ y $\int_{\gamma_3}\omega_2$. Comentar los resultados.
\end{enumerate}

\begin{indication}
En un segmento horizontal $\mathrm{d}y=0$ y en uno vertical $\mathrm{d}x=0$:
cada integral se reduce a una integral simple.
\end{indication}

\begin{solution}
\textbf{Pregunta 1.} Sirve $f(x,y)=xy$, pues $\partial f/\partial x=y$ y $\partial f/\partial y=x$.

\textbf{Pregunta 2.} Para una curva parametrizada por $t\mapsto\bigl(x(t),y(t)\bigr)$,
\[\int_\gamma\omega_1=\int\left[y(t)x'(t)+x(t)y'(t)\right]\,\mathrm{d}t.\]
El camino $\gamma_1$ es la unión de dos segmentos. En el primero, $x(t)=t$,
$y(t)=0$ y $t\in[0,1]$; por tanto, $\mathrm{d}x=\mathrm{d}t$ y
$\mathrm{d}y=0$. En el segundo, $x(t)=1$, $y(t)=t$, también con $t\in[0,1]$;
por tanto, $\mathrm{d}x=0$ y $\mathrm{d}y=\mathrm{d}t$. En consecuencia,
\[
\begin{aligned}
I_1=\int_{\gamma_1}\omega_1
&=\int_0^1(0\times1+t\times0)\,\mathrm{d}t+\int_0^1(t\times0+1\times1)\,\mathrm{d}t\\
&=0+\int_0^1\mathrm{d}t=1.
\end{aligned}
\]
Para $\gamma_2$, el primer segmento se parametriza mediante $x(t)=0$,
$y(t)=t$, y el segundo mediante $x(t)=t$, $y(t)=1$, con $t\in[0,1]$ en ambos:
\[
\begin{aligned}
I_2=\int_{\gamma_2}\omega_1
&=\int_0^1(t\times0+0\times1)\,\mathrm{d}t+\int_0^1(1\times1+t\times0)\,\mathrm{d}t\\
&=0+\int_0^1\mathrm{d}t=1.
\end{aligned}
\]
Los dos valores coinciden. Se podía prever sin cálculo: como
$\omega_1=\mathrm{d}f$ con $f(x,y)=xy$, su integral solo depende de los extremos,
\[\int_\gamma\omega_1=f(M)-f(O)=f(1,1)-f(0,0)=1.\]

\textbf{Pregunta 3.} Escribamos $\omega_2=A(x,y)\,\mathrm{d}x+B(x,y)\,\mathrm{d}y$.
Aquí $A(x,y)=y$ y $B(x,y)=0$. Si la forma fuera exacta, el criterio de Schwarz daría
\[\frac{\partial A}{\partial y}=\frac{\partial B}{\partial x}.\]
Sin embargo, $\partial A/\partial y=1$, mientras que $\partial B/\partial x=0$.
Las derivadas cruzadas difieren, por lo que $\omega_2$ no es exacta. Su integral
puede depender del camino seguido entre $O$ y $M$, como confirma el cálculo siguiente.

\textbf{Pregunta 4.} Como $\omega_2=y\,\mathrm{d}x$, solo puede contribuir a
la integral una variación de $x$ efectuada con un valor no nulo de $y$. Con las
parametrizaciones anteriores, sobre $\gamma_1$ se obtiene
\[
\begin{aligned}J_1=\int_{\gamma_1}\omega_2
&=\int_0^1 0\times1\,\mathrm{d}t+\int_0^1t\times0\,\mathrm{d}t\\&=0.\end{aligned}
\]
En $\gamma_2$, el primer segmento cumple $x(t)=0$, luego $\mathrm{d}x=0$; en
el segundo, $x(t)=t$, $y(t)=1$ y $\mathrm{d}x=\mathrm{d}t$. Así,
\[
\begin{aligned}J_2=\int_{\gamma_2}\omega_2
&=\int_0^1 t\times0\,\mathrm{d}t+\int_0^1 1\times1\,\mathrm{d}t\\&=1.\end{aligned}
\]
Por último, el segmento diagonal $\gamma_3$ se parametriza mediante
\[x(t)=t,\qquad y(t)=t,\qquad t\in[0,1],\]
de donde $\mathrm{d}x=\mathrm{d}t$ y $\mathrm{d}y=\mathrm{d}t$. Se sigue que
\[J_3=\int_{\gamma_3}\omega_2=\int_0^1 y(t)x'(t)\,\mathrm{d}t=\int_0^1 t\,\mathrm{d}t=\frac12.\]
Los tres caminos tienen los mismos extremos, pero dan tres valores distintos:
$J_1=0$, $J_2=1$ y $J_3=1/2$. Esta es precisamente la dependencia del camino
característica de una forma diferencial no exacta.
\end{solution}
\end{exo}

\begin{exo}[Misma variación de energía interna, tres modos de transferencia]{meme-delta-u-trois-transferts}
\keywords{primer principio, energía interna, calor, trabajo mecánico, trabajo eléctrico, calorimetría, función de estado}

Un líquido, su calorímetro y una pequeña resistencia sumergida constituyen un
sistema cerrado, macroscópicamente en reposo. El recipiente es rígido y se
supone constante la capacidad calorífica total del sistema:
\[C=2{,}00\ \mathrm{kJ\,K^{-1}}.\]
Se quiere llevar el sistema del estado de equilibrio $A$, a
$T_A=20{,}0\,^\circ\mathrm{C}$, al estado de equilibrio $B$, a
$T_B=25{,}0\,^\circ\mathrm{C}$. Se admite que
\[U(B)-U(A)=C(T_B-T_A).\]
Las variaciones de energía cinética y potencial macroscópicas son nulas. Se
desprecian todas las pérdidas hacia el exterior.

Se realizan independientemente los tres protocolos siguientes, partiendo del
mismo estado inicial $A$:
\begin{itemize}
\item[Protocolo 1] El sistema se pone en contacto térmico con una fuente caliente, sin suministrarle trabajo alguno.
\item[Protocolo 2] Las paredes están aisladas térmicamente. Una paleta agita el líquido gracias a la caída de una masa $m$ desde una altura $h=5{,}00$ m. La paleta y el líquido terminan en reposo, y toda la energía mecánica perdida por la masa se transfiere al sistema. Se tomará $g=9{,}81\ \mathrm{m\,s^{-2}}$.
\item[Protocolo 3] El sistema recibe un calor $Q_3=4{,}00$ kJ y la energía restante se aporta eléctricamente a la resistencia. La potencia eléctrica recibida es constante e igual a $\mathcal P=300$ W.
\end{itemize}

\begin{enumerate}
\item Calcular la variación de energía interna $\Delta U=U(B)-U(A)$ común a los tres protocolos.
\item Para cada uno de los dos primeros protocolos, determinar el calor $Q$ y el trabajo $W$ recibidos por el sistema. En el segundo, calcular la masa $m$ necesaria.
\item Para el tercer protocolo, calcular el trabajo eléctrico recibido y el tiempo de alimentación de la resistencia.
\item Las tres transformaciones tienen los mismos estados inicial y final. Explicar por qué sus valores de $Q$ y $W$ pueden ser, sin embargo, diferentes. ¿Contiene el sistema «calor» o «trabajo» en el estado final $B$?
\end{enumerate}

\begin{indication}
Aplicar $\Delta U=Q+W$ con el convenio del banquero. Para la paleta, el trabajo
recibido es igual a la disminución $mgh$ de la energía potencial de la masa.
La energía eléctrica que atraviesa la frontera por los cables se contabiliza
como trabajo eléctrico.
\end{indication}

\begin{solution}
\textbf{Pregunta 1.} La diferencia de temperatura vale $T_B-T_A=5{,}0$ K. Por tanto,
\[\Delta U=C(T_B-T_A)=2{,}00\times5{,}0=10{,}0\ \mathrm{kJ}.\]
Este valor solo depende de los dos estados de equilibrio $A$ y $B$.

\textbf{Pregunta 2.} En el primer protocolo no se recibe trabajo: $W_1=0$.
El primer principio da entonces
\[Q_1=\Delta U=10{,}0\ \mathrm{kJ}.\]
El calor es positivo porque la energía entra en el sistema.

En el segundo protocolo, las paredes están aisladas térmicamente, de modo que
$Q_2=0$. Toda la variación de energía interna procede del trabajo mecánico de la paleta:
\[W_2=\Delta U=10{,}0\ \mathrm{kJ}.\]
Como $W_2=mgh$,
\[m=\frac{W_2}{gh}=\frac{10{,}0\times10^3}{9{,}81\times5{,}00}\simeq204\ \mathrm{kg}.\]
El elevado orden de magnitud recuerda que incluso un aumento moderado de
temperatura corresponde ya a una energía mecánica importante.

\textbf{Pregunta 3.} El primer principio impone
\[W_3=\Delta U-Q_3=10{,}0-4{,}00=6{,}00\ \mathrm{kJ}.\]
Este trabajo es positivo: la alimentación eléctrica aporta energía al sistema.
Con $W_3=\mathcal P\,\Delta t$,
\[\Delta t=\frac{W_3}{\mathcal P}=\frac{6{,}00\times10^3}{300}=20{,}0\ \mathrm{s}.\]

\textbf{Pregunta 4.} Los tres caminos dan, respectivamente,
\[(Q_1,W_1)=(10{,}0,0)\ \mathrm{kJ},\qquad(Q_2,W_2)=(0,10{,}0)\ \mathrm{kJ},\]
y
\[(Q_3,W_3)=(4{,}00,6{,}00)\ \mathrm{kJ}.\]
En todos los casos, la suma $Q+W$ vale $10{,}0$ kJ y produce la misma variación
$\Delta U$. La energía interna es una función de estado, mientras que el calor
y el trabajo describen transferencias durante una transformación. En el estado
final $B$, el sistema posee energía interna, pero no contiene ni «calor» ni «trabajo».
\end{solution}
\end{exo}

\begin{exo}[Compresión bajo presión exterior constante]{compression-pression-exterieure-constante}
\seoready{true}
\keywords{trabajo de las fuerzas de presión, presión exterior constante, compresión, expansión, expansión libre}

Un gas se comprime desde un volumen $V_A$ hasta un volumen $V_B<V_A$ bajo
una presión exterior constante $P_0$.

\begin{enumerate}
\item Calcular el trabajo recibido por el gas.
\item El gas vuelve de $V_B$ a $V_A$ contra la misma presión exterior $P_0$. Calcular el trabajo recibido y compararlo con el de compresión.
\item Repetir la expansión de $V_B$ a $V_A$ cuando el gas se expande en el vacío. Interpretar los tres signos obtenidos.
\end{enumerate}

\begin{indication}
Utilizar $W=-\int P_{\mathrm{ext}}\,\mathrm{d}V$. Hay que integrar la presión
\emph{exterior}, incluso cuando la transformación no es cuasiestática.
\end{indication}

\begin{solution}
\textbf{Pregunta 1.} La presión exterior es constante, por lo que
\[W_{A\to B}=-P_0(V_B-V_A)=P_0(V_A-V_B)>0.\]
El gas recibe trabajo durante la compresión.

\textbf{Pregunta 2.} Para el retorno,
\[W_{B\to A}=-P_0(V_A-V_B)<0.\]
Los dos trabajos son opuestos porque las transformaciones se realizan contra
la misma presión exterior constante.

\textbf{Pregunta 3.} En el vacío, $P_{\mathrm{ext}}=0$, luego
\[W_{B\to A}^{\mathrm{vide}}=0.\]
Por tanto, una expansión no implica automáticamente un trabajo negativo: es
necesario que se rechace efectivamente una fuerza exterior.
\end{solution}
\end{exo}
